
\documentclass[11pt]{article}
\usepackage{series}

\renewcommand\SeriesShort{Delta Projection}
\renewcommand\PartRoman{XI}
\renewcommand\PartArabic{11}
\renewcommand\PartsInSeries{XI}

\title{Gaussian-Filtered Contact Loops and Finite One-Loop Tadpoles\\
\large Exact Euclidean Benchmark and the Local-Limit Counterterm}
\author{Peter Nero}
\date{July 2026, Version 1}

\begin{document}
\maketitle

\begin{abstract}
Earlier papers in this sequence established a structural dictionary in which Dirac delta
distributions arise as singular limits of finite kernels.  The present paper studies one exact
Euclidean benchmark obtained after declaring an external Gaussian momentum filter.  It does
not derive that external filter from an internal MTT fixed-point operator.

We study the Euclidean scalar one-loop tadpole
\[
I_0^{(d)}(m)=\int_{\mathbb R^d}\frac{d^d k}{(2\pi)^d}\frac{1}{k^2+m^2},
\]
which is ultraviolet divergent for \(d\ge 2\), and replace the point-contact propagator by the declared Gaussian-filtered propagator
\[
\Delta_\tau(k)=\frac{e^{-\tau k^2}}{k^2+m^2},\qquad \tau>0.
\]
We prove that
\[
I_\tau^{(d)}(m)=\int_{\mathbb R^d}\frac{d^d k}{(2\pi)^d}\frac{e^{-\tau k^2}}{k^2+m^2}
\]
is finite for every \(d\), \(m>0\), and \(\tau>0\), and derive the exact Schwinger representation
\[
I_\tau^{(d)}(m)
=(4\pi)^{-d/2}\int_0^\infty e^{-m^2 s}(s+\tau)^{-d/2}\,ds.
\]
In four dimensions this gives the closed form
\[
I_\tau^{(4)}(m)
=\frac{1}{16\pi^2}\left[
\frac{1}{\tau}
-m^2 e^{m^2\tau}E_1(m^2\tau)
\right],
\]
with the local-limit asymptotics
\[
I_\tau^{(4)}(m)
=
\frac{1}{16\pi^2}
\left[
\frac{1}{\tau}
+
m^2\bigl(\gamma+\log(m^2\tau)\bigr)
+O(m^4\tau|\log(m^2\tau)|)
\right].
\]
Thus the declared filter scale \(\tau\) gives a definite finite correction and an effective
Euclidean ultraviolet scale \(\Lambda_{\rm eff}\sim \tau^{-1/2}\).  The divergent local theory
is recovered as \(\tau\downarrow0\).  This proves neither all-loop finiteness nor physical
Lorentzian consistency.  An MTT interpretation additionally requires a selected
internal-to-external source map, while gauge/BRST identities, reflection positivity,
causality, and unitarity require separate certificates.
\end{abstract}

\section*{Version 1 Revision Note}
\begin{description}
\item[Supersedes] The unversioned April 2026 manuscript.
\item[Reason] The exact integral was correct, but the earlier wording treated the external
factor \(e^{-\tau k^2}\) as already selected by MTT fixed-point data.
\item[Resolution] The calculation is now stated as a declared-filter Euclidean benchmark.
The missing source intertwiner and the Lorentzian, gauge, and all-loop obligations are explicit.
\item[Retained result] The finiteness theorem, Schwinger representation, four-dimensional
closed form, and local-limit asymptotics are unchanged.
\item[Remaining boundary] MTT source promotion and physical QFT compatibility remain open.
\end{description}

\tableofcontents

\section{Purpose and claim discipline}

The previous structural papers showed that contact interactions hide products of delta kernels. For instance,
\[
\frac{\lambda}{4!}\int \phi(x)^4\,dx
\]
can be represented as a zero-width overlap vertex where all four field arguments are forced
to coincide.  The purpose of the present paper is to compute the first loop-level consequence
of a declared Gaussian external filter and to state exactly what would be required to promote
it to an MTT-selected physical kernel.

\begin{center}
\fbox{\parbox{0.9\linewidth}{
\textbf{Main claim.} In the Euclidean scalar contact-loop model, replacing the local
propagator by the declared Gaussian-filtered propagator
\[
\Delta_\tau(k)=e^{-\tau k^2}(k^2+m^2)^{-1}
\]
makes the one-loop tadpole finite for every \(\tau>0\). The ordinary ultraviolet divergence is recovered exactly as the local limit \(\tau\downarrow0\).
}}
\end{center}

\subsection{Non-claims}

This paper does not claim:
\begin{enumerate}
\item that a scalar Euclidean tadpole is a complete physical quantum field theory;
\item that arbitrary Gaussian smearing preserves Lorentzian unitarity, microcausality, or gauge/BRST identities;
\item that renormalization is unnecessary in the local limit;
\item that the numerical value of \(\tau\) is derived here from a specific Standard Model sector;
\item that one-loop finiteness proves all-loop finiteness or ultraviolet completion;
\item that the internal fixed-point gap by itself selects external momentum damping;
\item that this scalar calculation proves full MTT quantum gravity or full Standard Model phenomenology.
\end{enumerate}

The narrower result is the first execution-level loop theorem:
\[
\text{declared Gaussian filter}\Rightarrow\text{finite contact loop}\Rightarrow\text{controlled local-limit divergence}.
\]

\section{Declared proper-time filtering}

Let the Euclidean free scalar propagator be
\[
\Delta_0(k)=\frac{1}{k^2+m^2}.
\]
In flat Euclidean momentum variables, declare the model filter
\[
B_\tau(k)=e^{-\tau k^2/2},
\]
so that an internal line receives the factor
\[
B_\tau(k)\Delta_0(k)B_\tau(k)
=
\frac{e^{-\tau k^2}}{k^2+m^2}.
\]

\begin{definition}[Gaussian-filtered scalar propagator]
For \(m>0\) and \(\tau>0\), define
\[
\Delta_\tau(k):=\frac{e^{-\tau k^2}}{k^2+m^2}.
\]
The corresponding one-loop tadpole integral in \(d\) Euclidean dimensions is
\[
I_\tau^{(d)}(m):=\int_{\mathbb R^d}\frac{d^d k}{(2\pi)^d}\Delta_\tau(k).
\]
\end{definition}

\begin{remark}[Source boundary]
The factor \(e^{-\tau k^2}\) is an explicit external-model assumption in this paper.  An
internal heat kernel or fixed-point gap does not by itself imply damping in external momentum.
Promotion to an MTT-selected filter requires an intertwiner identifying the operator domains
and preserving the relevant symmetries.  The present paper proves only the scalar Euclidean
loop statement conditional on the displayed filter.
\end{remark}

\section{Finiteness theorem}

\begin{theorem}[Finite Gaussian-filtered tadpole]
Let \(d\ge1\), \(m>0\), and \(\tau>0\). Then
\[
I_\tau^{(d)}(m)
=
\int_{\mathbb R^d}\frac{d^d k}{(2\pi)^d}
\frac{e^{-\tau k^2}}{k^2+m^2}
\]
is finite.
\end{theorem}

\begin{proof}
For \(|k|\le1\), the integrand is bounded by \(m^{-2}\), hence the integral over the unit ball is finite. For \(|k|>1\),
\[
0\le \frac{e^{-\tau k^2}}{k^2+m^2}\le e^{-\tau k^2}.
\]
The Gaussian \(e^{-\tau k^2}\) is integrable over \(\mathbb R^d\) for every \(\tau>0\). Therefore the integral is finite.
\end{proof}

\begin{corollary}[Finite coincident propagator]
The Euclidean coincident kernel
\[
\Delta_\tau(x,x)
=
\int_{\mathbb R^d}\frac{d^d k}{(2\pi)^d}
\frac{e^{-\tau k^2}}{k^2+m^2}
\]
is finite for every \(d\), \(m>0\), and \(\tau>0\).
\end{corollary}

\section{Schwinger representation}

The following computation uses the standard proper-time method
\cite{Schwinger1951}.

\begin{theorem}[Exact proper-time representation]
For \(d\ge1\), \(m>0\), and \(\tau>0\),
\[
I_\tau^{(d)}(m)
=
(4\pi)^{-d/2}
\int_0^\infty e^{-m^2s}(s+\tau)^{-d/2}\,ds.
\]
\end{theorem}

\begin{proof}
Use the Schwinger representation
\[
\frac{1}{k^2+m^2}=\int_0^\infty e^{-s(k^2+m^2)}\,ds.
\]
Then
\[
I_\tau^{(d)}(m)
=
\int_0^\infty e^{-m^2s}
\left[
\int_{\mathbb R^d}\frac{d^d k}{(2\pi)^d}e^{-(s+\tau)k^2}
\right]ds.
\]
The inner Gaussian integral is
\[
\int_{\mathbb R^d}\frac{d^d k}{(2\pi)^d}e^{-(s+\tau)k^2}
=
(4\pi)^{-d/2}(s+\tau)^{-d/2}.
\]
Substitution gives the result. Tonelli's theorem justifies the interchange because the integrand is nonnegative.
\end{proof}

\begin{remark}
This formula shows exactly how the declared proper-time filtering changes the local loop.  In
the unfiltered case, the factor \(s^{-d/2}\) produces a short-proper-time divergence at
\(s=0\).  The filtered case shifts it to \((s+\tau)^{-d/2}\), removing the singularity for
every \(\tau>0\).
\end{remark}

\section{Four-dimensional tadpole}

Let \(d=4\). Then
\[
I_\tau^{(4)}(m)
=
\frac{1}{16\pi^2}
\int_0^\infty e^{-m^2s}(s+\tau)^{-2}\,ds.
\]

Let
\[
E_1(z):=\int_z^\infty \frac{e^{-u}}{u}\,du
\]
be the standard exponential integral.

\begin{theorem}[Closed form in four dimensions]
For \(m>0\) and \(\tau>0\),
\[
I_\tau^{(4)}(m)
=
\frac{1}{16\pi^2}
\left[
\frac{1}{\tau}
-
m^2e^{m^2\tau}E_1(m^2\tau)
\right].
\]
\end{theorem}

\begin{proof}
Set \(a=m^2\). We need
\[
J(a,\tau):=\int_0^\infty e^{-as}(s+\tau)^{-2}\,ds.
\]
With \(u=s+\tau\),
\[
J(a,\tau)=e^{a\tau}\int_\tau^\infty e^{-au}u^{-2}\,du.
\]
Integrating by parts, or differentiating \(E_1(a\tau)\), gives
\[
\int_\tau^\infty e^{-au}u^{-2}\,du
=
\frac{e^{-a\tau}}{\tau}
-
aE_1(a\tau).
\]
Thus
\[
J(a,\tau)=\frac{1}{\tau}-ae^{a\tau}E_1(a\tau).
\]
Multiplying by \(1/(16\pi^2)\) proves the formula.
\end{proof}

\begin{corollary}[Local-limit asymptotics in \(d=4\)]
As \(\tau\downarrow0\),
\[
I_\tau^{(4)}(m)
=
\frac{1}{16\pi^2}
\left[
\frac{1}{\tau}
+
m^2\bigl(\gamma+\log(m^2\tau)\bigr)
+
O\!\left(m^4\tau|\log(m^2\tau)|\right)
\right],
\]
where \(\gamma\) is Euler's constant.
\end{corollary}

\begin{proof}
For small \(z>0\),
\[
E_1(z)=-\gamma-\log z+z+O(z^2),
\qquad
e^z=1+O(z).
\]
Therefore
\[
-m^2e^{m^2\tau}E_1(m^2\tau)
=
m^2\bigl(\gamma+\log(m^2\tau)\bigr)
+
O(m^4\tau|\log(m^2\tau)|).
\]
Substitution into the closed form gives the result.
\end{proof}

\begin{remark}
The leading \(1/\tau\) term is the proper-time version of the usual quadratic divergence. The logarithmic term is the remnant of the usual logarithmic mass dependence. The local theory is recovered only as the singular limit \(\tau\downarrow0\).
\end{remark}

\section{General dimensional behavior}

The Schwinger formula gives the dimension dependence immediately.

\begin{proposition}[Local-limit divergence by dimension]
For fixed \(m>0\):
\begin{enumerate}
\item If \(d<2\), then \(I_\tau^{(d)}(m)\) has a finite limit as \(\tau\downarrow0\).
\item If \(d=2\), then \(I_\tau^{(2)}(m)\) diverges logarithmically.
\item If \(d>2\), then \(I_\tau^{(d)}(m)\) diverges as a positive power of \(1/\tau\), with leading order proportional to \(\tau^{1-d/2}\).
\end{enumerate}
\end{proposition}

\begin{proof}
The possible divergence comes only from the \(s=0\) region of
\[
\int_0^\infty e^{-m^2s}(s+\tau)^{-d/2}\,ds.
\]
For small \(s\), \(e^{-m^2s}=1+O(s)\). Thus the leading singularity is governed by
\[
\int_0^1 (s+\tau)^{-d/2}\,ds.
\]
This integral remains finite as \(\tau\downarrow0\) for \(d<2\), diverges logarithmically for \(d=2\), and diverges like \(\tau^{1-d/2}\) for \(d>2\).
\end{proof}

\section{One-loop mass correction in \texorpdfstring{\(\lambda\phi^4\)}{lambda phi4}}

In Euclidean \(\lambda\phi^4\) theory, the one-loop tadpole contribution to the two-point function has the schematic mass correction
\[
\delta m_\tau^2
=
\frac{\lambda}{2} I_\tau^{(d)}(m),
\]
up to convention-dependent factors.

In \(d=4\),
\[
\delta m_\tau^2
=
\frac{\lambda}{32\pi^2}
\left[
\frac{1}{\tau}
-
m^2e^{m^2\tau}E_1(m^2\tau)
\right].
\]

Thus the declared proper-time scale gives a definite finite correction:
\[
\boxed{
\delta m_\tau^2
\sim
\frac{\lambda}{32\pi^2}\frac{1}{\tau}
\quad
\text{for small }\tau.
}
\]

Equivalently, with
\[
\Lambda_{\rm eff}:=\tau^{-1/2},
\]
the leading correction scales as
\[
\delta m_\tau^2
\sim
\frac{\lambda}{32\pi^2}\Lambda_{\rm eff}^2.
\]

\begin{remark}
This is not yet a numerical prediction.  The physical value and external-domain meaning of
\(\tau\) require a source theorem or a fit to data.  Conditional on \(\tau\), the finite
correction is fixed.
\end{remark}

\section{Renormalization as the local-limit subtraction}

The filtered theory is finite for \(\tau>0\), but the local limit \(\tau\downarrow0\) is singular. In \(d=4\), the divergent pieces are
\[
\frac{1}{16\pi^2\tau}
+
\frac{m^2}{16\pi^2}\log(m^2\tau).
\]

A local renormalization scheme subtracts the singular terms as \(\tau\downarrow0\), in line
with standard local constructions \cite{EpsteinGlaser1973}.  In the present interpretation,
this is not a contradiction of finite coherent kernels. It says:

\begin{center}
\fbox{\parbox{0.84\linewidth}{\centering
Renormalization is the repair required when finite coherent overlap is forced into the
zero-width local limit.}}
\end{center}

\begin{definition}[Local-limit counterterm]
A local-limit mass counterterm in \(d=4\) is a \(\tau\)-dependent term \(\delta m_{\rm ct}^2(\tau)\) chosen so that
\[
\delta m_{\rm ren}^2
=
\frac{\lambda}{2}I_\tau^{(4)}(m)
+
\delta m_{\rm ct}^2(\tau)
\]
has a finite limit or a prescribed finite value as \(\tau\downarrow0\).
\end{definition}

\begin{remark}
This paper does not replace standard renormalization. It identifies the finite coherent object whose local limit standard renormalization repairs.
\end{remark}

\section{Phenomenological scale and bounds}

At external momentum scale \(E\), the filtered propagator differs from the local propagator by
\[
\frac{\Delta_\tau(k)}{\Delta_0(k)}
=
e^{-\tau k^2}
=
1-\tau k^2+O(\tau^2 k^4).
\]
Therefore low-energy recovery requires
\[
\tau E^2\ll1,
\]
or
\[
\Lambda_{\rm eff}=\tau^{-1/2}\gg E.
\]

\begin{proposition}[Model-independent low-energy bound]
If experiments confirm the local propagator structure up to relative error \(\varepsilon\) at characteristic Euclidean momentum scale \(E\), then the scalar filtered model requires
\[
\tau E^2\lesssim \varepsilon,
\qquad
\text{equivalently}\qquad
\Lambda_{\rm eff}\gtrsim \frac{E}{\sqrt{\varepsilon}}.
\]
\end{proposition}

\begin{proof}
For \(z=\tau E^2\ll1\),
\[
|e^{-z}-1|\le z+O(z^2).
\]
Demanding this be \(\lesssim\varepsilon\) gives the stated bound.
\end{proof}

\begin{remark}
This is the first phenomenological form of the execution program. It does not yet identify a real experiment or sector-specific value of \(\tau\). It shows how \(\tau\) enters observables once the coherent kernel is fixed.
\end{remark}

\section{Scope of proof}

\begin{center}
\begin{tabular}{p{0.36\linewidth}p{0.54\linewidth}}
\hline
\textbf{Level} & \textbf{Status in this paper}\\
\hline
Scalar Euclidean tadpole & Proved finite for \(\tau>0\).\\
Schwinger representation & Proved exactly.\\
Four-dimensional asymptotics & Proved using \(E_1\).\\
Mass correction formula & Derived in the standard \(\lambda\phi^4\) tadpole convention.\\
Renormalization interpretation & Structural interpretation of the local limit.\\
Lorentzian unitarity & Not proved here.\\
Gauge/BRST preservation & Not proved here.\\
Numerical value of \(\tau\) & Not derived here.\\
MTT source of the external filter & Not proved here.\\
All-loop UV completion & Not implied by this one-loop result.\\
Full Standard Model prediction & Not claimed.\\
\hline
\end{tabular}
\end{center}

\section{Conclusion}

This paper gives an exact loop-level benchmark in the delta/projection program.  Starting
from the declared Gaussian external filter, the scalar propagator becomes
\[
\Delta_\tau(k)=\frac{e^{-\tau k^2}}{k^2+m^2}.
\]
The one-loop tadpole
\[
I_\tau^{(d)}(m)
=
\int\frac{d^d k}{(2\pi)^d}
\frac{e^{-\tau k^2}}{k^2+m^2}
\]
is finite for every \(d\), \(m>0\), and \(\tau>0\). In \(d=4\), it has the exact form
\[
I_\tau^{(4)}(m)
=
\frac{1}{16\pi^2}
\left[
\frac{1}{\tau}
-
m^2e^{m^2\tau}E_1(m^2\tau)
\right],
\]
and the singular local limit reproduces the usual ultraviolet divergence.

The exact benchmark result is:
\[
\boxed{
\text{a declared Gaussian external filter yields a specific finite one-loop correction.}
}
\]

The MTT task is first to derive the same external filter from selected geometry.  Only then do
multi-loop control, Lorentzian causal continuation, reflection positivity, and
gauge/BRST-compatible sectors become physical promotion questions.

\begin{thebibliography}{9}

\bibitem{Schwinger1951}
J. Schwinger,
``On Gauge Invariance and Vacuum Polarization,''
\emph{Physical Review} \textbf{82} (1951), 664--679.
\url{https://doi.org/10.1103/PhysRev.82.664}

\bibitem{EpsteinGlaser1973}
H. Epstein and V. Glaser,
``The role of locality in perturbation theory,''
\emph{Annales de l'Institut Henri Poincar\'e A} \textbf{19} (1973), 211--295.
\url{https://eudml.org/doc/75788}

\end{thebibliography}

\end{document}
